\section{Discussion}
\label{sec:discussion}

The central observation of this paper is that stereographic
projection, composed with standard QSP, converts the bounded
polynomial output of the circuit into an unbounded rational
function---without modifying the quantum algorithm itself.  The
basis change identity
$S_z = V^{-1}W(\rtilde)V$ with $V = \mathrm{diag}(1,i)$
(Proposition~\ref{prop:basis-change}) makes this reduction
precise: the full machinery of standard QSP---phase-finding
algorithms~\cite{DongMengWhaleyEtAl2021,Haah2019}, the
function class characterization, convergence theory---applies
immediately, and the ``unboundedness'' is a free consequence of
the encoding/decoding interface.  The stereographic viewpoint
also illuminates standard QSP itself: even the standard
protocol, when decoded via stereographic projection, produces
unbounded rational functions of the real-line parameter
(Section~\ref{subsec:further-consequences}).

At the same time, the construction has clear limitations.
Stereographic block encoding requires knowledge of the
eigenbasis (Section~\ref{sec:block-encoding}), restricting
applicability to structured Hamiltonians---we have given
explicit circuits for diagonal, Pauli-$Z$ sum, and Heisenberg
models, but these do not cover the general case addressed by
QSVT~\cite{GSLW2019}.  The measurement overhead grows as
$\mathcal{O}(r^6/\epsilon^2)$ for large eigenvalues
(Eq.~\eqref{eq:shot-count}), making the approach impractical
when $r \gg 1$.  The practical advantage window is
depth-limited devices processing operators with moderate
eigenvalues and known eigenbasis structure.

Our construction rests on the QSP/QSVT framework established
by Low and Chuang~\cite{LowChuang2017,LowChuang2019} and
Gily\'en et al.~\cite{GSLW2019}, with comprehensive reviews by
Martyn et al.~\cite{MartynRossiTanChuang2021} and
Lin~\cite{Lin2022}.  The formal separation between bounded and
unbounded polynomial approximations established by Tang and
Tian~\cite{TangTian2024} provides the theoretical motivation:
it shows that the boundedness constraint is a genuine barrier,
not merely a technical inconvenience.  Among extensions,
GQSP~\cite{MotlaghWiebe2023} is the closest in spirit: it
lifts the parity and realness constraints on~$P$ and provides
dramatically faster angle-finding, but retains the boundedness
constraint $|P| \leq 1$.  Stereographic QSP is
complementary---it lifts boundedness but not parity, and
produces rational rather than polynomial output.  Quantum Phase
Processing~\cite{WangZhangYuWang2023}, modular
QSP~\cite{RossiCeroniChuang2025}, Randomized
QSVT~\cite{WangZhangHazraEtAl2025}, Hamiltonian block
encoding~\cite{KangSu2025}, and GQSP-based matrix function
synthesis~\cite{MahasingheEtAl2025} each address different
costs of the standard framework while retaining bounded output;
our work changes the classical encoding/decoding interface
rather than the quantum algorithm.  The connection to Boyd's
rational Chebyshev
functions~\cite{Boyd1990,Boyd1987} links stereographic QSP to
classical approximation theory on unbounded domains, suggesting
that convergence rates for rational Chebyshev expansions
translate directly to QSP performance guarantees.

The framework opens several directions for further work.  On
the algorithmic side, finding QSP phases for a given target
rational function $f(r)$ requires decomposing $f$ into a
polynomial pair $(P, Q)$ satisfying the normalization
constraint before invoking standard phase-finding; GQSP's
efficient angle-finding~\cite{MotlaghWiebe2023} may simplify
this step.  On the theoretical side, the numerical results in
Figure~\ref{fig:target-gallery} demonstrate convergence for
several targets, with the inset in panel~(e) showing that
eigenvalue inversion converges fastest while bounded targets
converge more slowly---understanding the precise relationship
between rational Chebyshev convergence rates and the QSP
normalization constraint is an open question.  The natural
error metric on unbounded domains is the weighted $L^2$ norm
with Cauchy--Lorentz weight $1/(1+r^2)$, which downweights
large-$r$ errors where stereographic compression is most
severe.

Extending the class of tractable Hamiltonians---sparse models
with known eigenbasis structure, integrable models, or
tensor-network states---would establish the practical scope of
the approach, and the key question is whether the eigenbasis
requirement can be relaxed, perhaps by combining stereographic
encoding with quantum phase estimation.  Replacing the phase
rotations $\exp[i\phi\PauliZ]$ with general signal-type
operators $O_\alpha$ parameterized by $\alpha \in \hat{\C}$
yields a normalization condition
\begin{equation}
	|P|^2 + |Q|^2 = (1{+}r^2)^d \textstyle\prod_j(1{+}|\alpha_j|^2),
	\label{eq:generalized-norm}
\end{equation}
whose forward direction follows from the same inductive
argument as Theorem~\ref{thm:unnormalized}, but whose
converse---decomposing arbitrary polynomial pairs into such a
product---remains open.  Finally, extending stereographic
encoding to the full QSVT framework, where the signal is a
block-encoded matrix rather than a scalar, requires defining a
matrix-valued stereographic encoding and analyzing its
interaction with the QSVT product structure.

The stereographic projection is one of the oldest maps in
mathematics, but its application to quantum signal processing
appears to be new.  The fact that a simple change of
coordinates can bypass a fundamental constraint of QSP suggests
that the interface between encoding and processing---the way
classical data enters and exits the quantum circuit---deserves
systematic study.  Other geometric constructions (hyperbolic
encodings, projective coordinates) may yield analogous
extensions, and the general question of which function classes
are realizable under different encoding/decoding strategies
remains open.
